\section*{Experimental Design}

The experiment was a discrete-choice task involving repeated decisions between three transportation modes: bus, metro, and ride-hailing. Each alternative was described by four attributes: cost, time, comfort, and CO\textsubscript{2} emissions. The study comprised 150 trials per participant, divided into two phases. In the training phase, participants completed 100 trials in a stable context. The colors and positions of rows (attribute labels) and columns (option labels) were fixed during this phase. In the testing phase, participants completed 50 trials in which the position of attribute rows was shuffled every five trials. This manipulation was designed to disrupt attentional habits and re-engage exploratory, goal-directed attention \citep{dayan2000,gottlieb2012}.

Attribute values in each trial were sampled through a standard DCE design procedure intended to ensure that no alternative dominated the others. In other words, no alternative was objectively better than the others across all attributes. The key manipulation was not the economic content of the options, but the stability of the visual environment. The stable phase allowed participants to learn the display. The disrupted phase broke that stable environment.

\begin{figure}[H]
    \centering
    \includegraphics[width=0.78\linewidth]{word_media/media/image1.png}
    \caption{Choice interface used in the discrete-choice task. Participants repeatedly chose between three transportation options described by cost, time, comfort, and CO\textsubscript{2} emissions.}
    \label{fig:framework}
\end{figure}
